\chapter{Convertibles and the Credit-Equity Link}\label{ch:dv:convertibles-and-the-credit-equity-link}

A company whose shares trade at 30 issues a five-year bond that the holder may exchange for 2.5 shares at any time. With the
share at 80 the bond is worth 202, barely more than the 200 its shares would fetch: it trades like the share. With the share at
10 it is worth 69, barely more than the 68 of a plain bond of the same company, and that plain bond has itself fallen, because a
company whose share has collapsed is more likely to default. It trades on its credit. In between the convertible is both at once,
and a convertible desk lives in that region. It holds the bond, sells the share against it, buys credit protection, and earns the
convexity. This chapter describes the instrument and its \omterm{def:dv:convertibles-and-the-credit-equity-link:floor}{bond floor}, \omterm{def:dv:convertibles-and-the-credit-equity-link:value}{conversion value} and convex middle. It then builds an
\omterm{def:dv:convertibles-and-the-credit-equity-link:e2c}{equity-to-credit model} on a finite-difference grid, derives what the desk hedges, and ends with what happened when the arbitrage's
financing disappeared.

\section{The instrument}

\begin{definition}[Convertible bond]\label{def:dv:convertibles-and-the-credit-equity-link:cb}
A \emph{convertible bond}\index{convertible bond} is a bond that its holder may exchange for a fixed number of the issuer's shares,
at any time or during set periods, instead of receiving coupons and principal; it usually also gives the issuer the right to call
it and sometimes the holder the right to put it.
\end{definition}

\begin{definition}[Conversion ratio]\label{def:dv:convertibles-and-the-credit-equity-link:ratio}
The \emph{conversion ratio}\index{conversion ratio} is the number of shares received per bond on conversion.
\end{definition}

\begin{definition}[Conversion price]\label{def:dv:convertibles-and-the-credit-equity-link:price}
The \emph{conversion price}\index{conversion price} is the face value divided by the \omterm{def:dv:convertibles-and-the-credit-equity-link:ratio}{conversion ratio}: the share price at which
converting is worth the face value.
\end{definition}

\begin{definition}[Conversion value]\label{def:dv:convertibles-and-the-credit-equity-link:value}
The \emph{conversion value}\index{conversion value} (parity) is the \omterm{def:dv:convertibles-and-the-credit-equity-link:ratio}{conversion ratio} times the current share price, what the bond
would be worth if converted now.
\end{definition}

\begin{definition}[Conversion premium]\label{def:dv:convertibles-and-the-credit-equity-link:premium}
The \emph{conversion premium}\index{conversion premium} is the convertible's price over its \omterm{def:dv:convertibles-and-the-credit-equity-link:value}{conversion value}, minus one.
\end{definition}

The chapter's bond has a face of 100, five years to maturity, an annual coupon of 2\%, and a \omterm{def:dv:convertibles-and-the-credit-equity-link:ratio}{conversion ratio} of 2.5, so a \omterm{def:dv:convertibles-and-the-credit-equity-link:price}{conversion
price} of 40 against a share at 30. The issuer may call it at 100 from the end of year two, but only if the share is at or above 52,
130\% of the \omterm{def:dv:convertibles-and-the-credit-equity-link:price}{conversion price}. Rates are 3\%, the dividend yield 1\% and the share's volatility 30\%. The issuer's credit spread is 300
basis points, a default intensity of 5\% a year with 40\% recovery of face.

\begin{definition}[Soft call]\label{def:dv:convertibles-and-the-credit-equity-link:softcall}
A \emph{soft call}\index{soft call} is an issuer's right to redeem a convertible early that can be exercised only if the share
price has traded above a trigger level (for instance 130\% of the \omterm{def:dv:convertibles-and-the-credit-equity-link:price}{conversion price}): a call forces holders
to convert when conversion is worth more than the call price.
\end{definition}

The issuer's call is the holder's cost. Calling forces conversion when the \omterm{def:dv:convertibles-and-the-credit-equity-link:value}{conversion value} exceeds the call price. That caps the
bond's value near the trigger and removes the time value the holder would otherwise keep. In the model below the \omterm{def:dv:convertibles-and-the-credit-equity-link:softcall}{soft call} is worth
3.9 at a share price of 40 and 4.7 at 52: 115.3 and 138.6 with the call, against 119.2 and 143.3 without it.

\section{Bond floor, conversion value and the convex region}

\begin{definition}[Bond floor]\label{def:dv:convertibles-and-the-credit-equity-link:floor}
The \emph{bond floor}\index{bond floor} of a convertible is the value of the same bond without the conversion right: its coupons and
principal discounted for the issuer's credit (Book 2, chapter 21).
\end{definition}

A convertible is worth at least its \omterm{def:dv:convertibles-and-the-credit-equity-link:floor}{bond floor}, because the holder can simply keep the bond, and at least its \omterm{def:dv:convertibles-and-the-credit-equity-link:value}{conversion value},
because it can convert. It is worth more than both, by the value of the option to choose later. Far from the \omterm{def:dv:convertibles-and-the-credit-equity-link:price}{conversion price} one
bound dominates and the convertible behaves like it: a bond on the left, the shares on the right. Near the \omterm{def:dv:convertibles-and-the-credit-equity-link:price}{conversion price} both
matter. There the convertible is convex in the share, and that convexity is what the arbitrage buys
(\cref{fig:dv:convertibles-and-the-credit-equity-link:profile}).

\begin{example}[The convertible across share prices]\label{ex:dv:convertibles-and-the-credit-equity-link:profile}
Under a constant 5\% hazard, the convertible is worth 89.1, 100.2, 116.0, 156.7 and 203.2 at share prices of 20, 30, 40, 60 and 80. Its \omterm{def:dv:convertibles-and-the-credit-equity-link:floor}{bond
floor} is 83.2 throughout. The \omterm{def:dv:convertibles-and-the-credit-equity-link:value}{conversion values} are 50, 75, 100, 150 and 200, so the \omterm{def:dv:convertibles-and-the-credit-equity-link:premium}{conversion premium} falls from 78\% at 20 to 34\% at 30,
16\% at 40, 4.5\% at 60 and 1.6\% at 80. The delta rises from 0.84 shares per bond at 20 to 2.39 at 80, approaching the \omterm{def:dv:convertibles-and-the-credit-equity-link:ratio}{conversion ratio} of 2.5.
\end{example}

\begin{omfigure}
\begin{tikzpicture}
\begin{axis}[width=10.5cm,height=5.4cm,
  xlabel={share price},ylabel={value per 100 of face},
  tick label style={font=\small},label style={font=\small},
  xmin=10,xmax=90,ymin=0,ymax=230,grid=major,grid style={gray!25},
  legend columns=3,legend cell align=left,legend style={font=\footnotesize,at={(0.5,-0.28)},anchor=north,draw=none,fill=none,/tikz/every even column/.append style={column sep=6pt}},
  table/col sep=comma]
\addplot[omThm,very thick] table[x=s,y=cb]{figdata/derivatives/21-convertibles-and-the-credit-equity-link/profile.csv};
\addplot[omThm,thick,dashed] table[x=s,y=cb_const]{figdata/derivatives/21-convertibles-and-the-credit-equity-link/profile.csv};
\addplot[omDef,very thick] table[x=s,y=floor]{figdata/derivatives/21-convertibles-and-the-credit-equity-link/profile.csv};
\addplot[omDef,thick,dashed] table[x=s,y=floor_const]{figdata/derivatives/21-convertibles-and-the-credit-equity-link/profile.csv};
\addplot[gray,thick,dotted] table[x=s,y=conv]{figdata/derivatives/21-convertibles-and-the-credit-equity-link/profile.csv};
\legend{convertible (equity-to-credit),convertible (constant hazard),bond floor (equity-to-credit),bond floor (constant hazard),conversion value}
\end{axis}
\end{tikzpicture}

\omcaption{The chapter's convertible (\omterm{def:dv:convertibles-and-the-credit-equity-link:price}{conversion price} 40, share at 30 today) by share price, under a constant hazard and under an
equity-to-credit hazard that rises as the share falls. The \omterm{def:dv:convertibles-and-the-credit-equity-link:floor}{bond floor} is flat in the first model and falls with the share in the
second; the convertible follows the \omterm{def:dv:convertibles-and-the-credit-equity-link:value}{conversion value} on the right and the floor on the left. Data: the
tutorial.}\label{fig:dv:convertibles-and-the-credit-equity-link:profile}
\end{omfigure}

\section{Equity-to-credit models}

The constant-hazard picture has a flaw that matters most where the desk's risk is largest. When a company's share falls far, its
credit usually deteriorates too, and the \omterm{def:dv:convertibles-and-the-credit-equity-link:floor}{bond floor}, on which the whole left side of the profile rests, falls with it. Models that tie
default to the share price capture that.

\begin{definition}[Equity-to-credit model]\label{def:dv:convertibles-and-the-credit-equity-link:e2c}
An \emph{equity-to-credit model}\index{equity-to-credit model} prices equity and credit instruments of one issuer together by making the
default intensity a decreasing function of the share price, for instance $\lambda(S)=\lambda_0(S/S_0)^{-p}$. At default the share
jumps to zero and bondholders recover a fraction of face.
\end{definition}

Under the pricing measure the share must still earn the risk-free rate on average, so its drift is raised by $\lambda(S)$ to compensate
for the jump to zero. The convertible's value $V(t,S)$ solves, in $x=\ln S$,
\[
\partial_tV+\tfrac12\sigma^2\partial_{xx}V+\bigl(r-q+\lambda(S)-\tfrac12\sigma^2\bigr)\partial_xV-\bigl(r+\lambda(S)\bigr)V+\lambda(S)\,RF=0,
\]
between coupon dates, with the constraints $V\ge\kappa S$ for conversion, $V\le\max(\text{call price},\kappa S)$ when the call is live, and
$V\ge\text{put price}$ on put dates. It is the \omterm{def:dv:black-scholes-three-ways:pde}{Black--Scholes equation} (chapter 3) with a killing rate and a recovery source. Ayache, Forsyth
and Vetzal set out this single-equation formulation. Tsiveriotis and Fernandes had earlier split the convertible into a cash-only part
discounted at the risky rate and an equity part discounted at the risk-free rate. The build solves the equation by Crank--Nicolson on a
500-point log-price grid (finite differences, One Quant Book 4, chapter 27, and chapter 22 of this book). It applies the constraints after
each step, the same projection that priced American options in chapter 6. Andersen and Buffum showed that naive calibration of such models
can bias prices significantly, and that the hazard function should be calibrated jointly to the issuer's credit curve and to its options.

\begin{omfigure}
\begin{tikzpicture}
\begin{axis}[width=10.5cm,height=4.6cm,
  xlabel={share price},ylabel={default intensity (\% a year)},
  tick label style={font=\small},label style={font=\small},
  xmin=10,xmax=90,ymin=0,ymax=20,grid=major,grid style={gray!25},
  legend columns=2,legend cell align=left,legend style={font=\footnotesize,at={(0.5,-0.32)},anchor=north,draw=none,fill=none,/tikz/every even column/.append style={column sep=8pt}},
  table/col sep=comma]
\addplot[omThm,very thick] table[x=s,y=hazard]{figdata/derivatives/21-convertibles-and-the-credit-equity-link/profile.csv};
\draw[omDef,thick,dashed] (axis cs:10,5) -- (axis cs:90,5);
\legend{{$\lambda(S)=5\%\,(S/30)^{-1.2}$},constant 5\%}
\end{axis}
\end{tikzpicture}

\omcaption{The two hazards of the chapter, equal at today's share price of 30 (a 300 basis-point spread with 40\% recovery). The
equity-to-credit hazard is 18.7\% a year at a share price of 10 and 1.5\% at 80; the dashed line is the constant hazard. Data:
the chapter's code.}\label{fig:dv:convertibles-and-the-credit-equity-link:hazard}
\end{omfigure}

\begin{example}[What the credit link changes]\label{ex:dv:convertibles-and-the-credit-equity-link:e2c}
Take $\lambda(S)=5\%\,(S/30)^{-1.2}$, the same 300-basis-point spread at today's price. The hazard is 18.7\% a year at a share price of 10
(\cref{fig:dv:convertibles-and-the-credit-equity-link:hazard}). The \omterm{def:dv:convertibles-and-the-credit-equity-link:floor}{bond
floor} is now 78.1 at 20 and 90.4 at 80, where it was 83.2 at both. The convertible is worth 83.9 at 20 against 89.1 under the constant hazard, and
98.5 at 30 against 100.2. The delta is 1.55 shares per bond at 30 against 1.37, and 1.39 against 0.84 at 20. A falling share now also hurts the \omterm{def:dv:convertibles-and-the-credit-equity-link:floor}{bond
floor}, so the convertible's value falls faster. Below a share price of about 19 the gamma turns negative: $-0.03$ at 15 and $-0.10$ at 10. The convexity the desk thought it owned disappears
where the credit takes over.
\end{example}

\begin{omfigure}
\begin{tikzpicture}
\begin{groupplot}[group style={group size=2 by 1,horizontal sep=1.6cm},
  width=6.6cm,height=5cm,
  tick label style={font=\small},label style={font=\small},grid=major,grid style={gray!25},table/col sep=comma]
\nextgroupplot[xlabel={share price},ylabel={delta (shares per bond)},xmin=10,xmax=90,ymin=0,ymax=2.6,
  title={\small delta},
  legend columns=1,legend cell align=left,legend style={font=\footnotesize,at={(0.5,-0.32)},anchor=north,draw=none,fill=none}]
\addplot[omThm,very thick] table[x=s,y=delta]{figdata/derivatives/21-convertibles-and-the-credit-equity-link/profile.csv};
\addplot[omDef,thick,dashed] table[x=s,y=delta_const]{figdata/derivatives/21-convertibles-and-the-credit-equity-link/profile.csv};
\legend{equity-to-credit,constant hazard}
\nextgroupplot[xlabel={share price},ylabel={gamma},xmin=10,xmax=90,ymin=-0.12,ymax=0.08,
  yticklabel style={/pgf/number format/fixed,/pgf/number format/precision=2},
  title={\small gamma}]
\addplot[omThm,very thick] table[x=s,y=gamma]{figdata/derivatives/21-convertibles-and-the-credit-equity-link/profile.csv};
\addplot[omDef,thick,dashed] table[x=s,y=gamma_const]{figdata/derivatives/21-convertibles-and-the-credit-equity-link/profile.csv};
\draw[gray] (axis cs:10,0) -- (axis cs:90,0);
\end{groupplot}
\end{tikzpicture}

\omcaption{The convertible's delta and gamma under the two models. With the credit link the delta stays high as the share falls, because the
\omterm{def:dv:convertibles-and-the-credit-equity-link:floor}{bond floor} falls with it, and the gamma turns negative below about 19: the long-convexity position the arbitrage relies on becomes short
convexity in distress. Data: the tutorial.}\label{fig:dv:convertibles-and-the-credit-equity-link:greeks}
\end{omfigure}

\section{What a convertible desk hedges}

The classic position is long the convertible and short delta shares. Hedged that way it profits from realised volatility above the implied
volatility in its price (chapter 4), as a long option does. It keeps the credit risk of the \omterm{def:dv:convertibles-and-the-credit-equity-link:floor}{bond floor} and the risk of the model's delta.
Desks hedge the credit with a credit default swap (Book 2, chapter 23) sized to the convertible's credit sensitivity, and they worry most
about the combined scenario.

\begin{definition}[Convertible arbitrage]\label{def:dv:convertibles-and-the-credit-equity-link:arb}
\emph{Convertible arbitrage}\index{convertible arbitrage} is the strategy of buying \omterm{def:dv:convertibles-and-the-credit-equity-link:cb}{convertible bonds} that trade cheap to a model's value,
selling the issuer's shares short in the model's delta and often buying credit protection, to earn the convertible's convexity, its coupon
net of the short's cost, and the convergence of its price to value.
\end{definition}

\begin{example}[The hedge book]\label{ex:dv:convertibles-and-the-credit-equity-link:hedge}
At a share price of 30 the \omterm{def:dv:convertibles-and-the-credit-equity-link:e2c}{equity-to-credit model}'s delta is 1.55 shares per bond. A 1 basis point widening of the issuer's spread, at a fixed share
price, costs the convertible 0.0167 per bond (1.67 per 100 basis points), about half the 0.0296 of the \omterm{def:dv:convertibles-and-the-credit-equity-link:floor}{bond floor}, since the conversion right
cushions it. With a five-year risky annuity of 4.12, credit protection on 40.6\% of the face offsets that sensitivity.
\end{example}

\begin{example}[Share down 20\%, spread 300 basis points wider]\label{ex:dv:convertibles-and-the-credit-equity-link:stress}
The share falls from 30 to 24, and the issuer's spread at that price ends 300 basis points above today's, a hazard of 10\%. The convertible falls from
98.47 to 86.62, a loss of 11.84 per 100 of face. The short of 1.55 shares gains 9.30. The delta-hedged position loses 2.55 per 100 of face, 2.55
million on 100 million. Hedged with the constant-hazard model's delta of 1.37 instead, it loses 4.38. The credit protection gains 4.47 in the same
scenario and turns the loss into a gain of 1.93. The wider spread alone, with no move in the share, costs the delta-hedged position 3.35. The share
move alone, with the spread held, earns it 2.21 of convexity (\cref{fig:dv:convertibles-and-the-credit-equity-link:stress}).
\end{example}

\begin{omfigure}
\begin{tikzpicture}
\begin{axis}[width=10.5cm,height=5cm,
  xlabel={share move (\%)},ylabel={delta-hedged P\&L (per 100 of face)},
  tick label style={font=\small},label style={font=\small},
  xmin=-32,xmax=22,ymin=-4,ymax=4.5,grid=major,grid style={gray!25},
  legend columns=3,legend cell align=left,legend style={font=\footnotesize,at={(0.5,-0.3)},anchor=north,draw=none,fill=none,/tikz/every even column/.append style={column sep=8pt}},
  table/col sep=comma]
\addplot[omThm,very thick,mark=*,mark size=1.4pt] table[x=move,y=w0]{figdata/derivatives/21-convertibles-and-the-credit-equity-link/stress.csv};
\addplot[omDef,very thick,mark=square*,mark size=1.4pt] table[x=move,y=w150]{figdata/derivatives/21-convertibles-and-the-credit-equity-link/stress.csv};
\addplot[omProp,very thick,mark=triangle*,mark size=1.6pt] table[x=move,y=w300]{figdata/derivatives/21-convertibles-and-the-credit-equity-link/stress.csv};
\draw[gray] (axis cs:-32,0) -- (axis cs:22,0);
\legend{spread unchanged,+150 bp,+300 bp}
\end{axis}
\end{tikzpicture}

\omcaption{The long convertible, short 1.55 shares per bond (the equity-to-credit delta at a share price of 30), after a move in the share and a
widening of the issuer's spread at the new share price. The convexity pays when the spread holds; a widening turns every scenario into a loss. Data:
the tutorial.}\label{fig:dv:convertibles-and-the-credit-equity-link:stress}
\end{omfigure}

The model's choice of delta is itself a hedge decision. An equity-to-credit delta is larger, because it counts the credit channel through which a
falling share lowers the \omterm{def:dv:convertibles-and-the-credit-equity-link:floor}{bond floor}, and so it protects more in the scenario that hurts. A desk running the constant-hazard delta is implicitly
long credit through its equity hedge. It is also exposed to the jump to default, which a delta and a CDS sized for small moves hedge only in part
(One Quant Book 6, chapter 13, and the structural models of chapter 14 there).

\section{Convertible arbitrage and its crises}

The arbitrage earns its return from the difference between a model's value and a market price, and from the convexity it holds. It is financed:
the fund borrows against the convertibles and lends the shares it sells through its prime broker. That financing is the arbitrage's weak point. When
many funds hold the same bonds and must sell at the same time, the prices fall together, away from any model's value. The hedges then protect
against the wrong risk.

Mitchell, Pedersen and Pulvino's ``Slow moving capital'' (2007) and Mitchell and Pulvino's study of 2008 (2012) examine what happens then.
In their account, hedge funds held up to 75\% of the convertible market; in early 2005 large investors began to withdraw capital from
convertible-arbitrage funds, more than 20\% of it in the first quarter (a figure they take from the Barclay Group), and the funds' forced
sales pushed convertibles below their fundamental values. For 2008, Mitchell and Pulvino describe the mechanism. The imminent failure of large prime brokers abruptly cut the leverage available to hedge funds,
seemingly long-term financing became short-term, and relative-value funds could no longer hold similar assets at similar prices. For a model this
means a scenario the stress example above does not contain: the convertible cheapens relative to its own model value while the share and credit
barely move. No hedge in the share or the CDS covers it. Its protection is lower leverage and financing that cannot be withdrawn overnight.

\section{Tutorial: a convertible on a grid}\label{tut:dv:convertibles-and-the-credit-equity-link:tut}

\begin{tutorial}
\textbf{Goal.} Price the chapter's convertible under a constant and an equity-to-credit hazard, compute its delta, gamma and credit sensitivity,
and stress the hedged position. \textbf{End state:} the four figures and the numbers of the weekend problem.
\begin{tutsteps}
\item \textbf{The operator}: drift raised by the hazard, killing at $r+\lambda$, recovery as a source:
\omcode{code/firm/convertible/firm_convertible.py}{70}{74}{The convertible's pricing operator.}\label{lst:dv:convertibles-and-the-credit-equity-link:operator}
\item \textbf{The constraints after each step}: coupons, puts, conversion and the \omterm{def:dv:convertibles-and-the-credit-equity-link:softcall}{soft call}:
\omcode{code/firm/convertible/firm_convertible.py}{92}{100}{Coupons, puts, conversion and the call.}\label{lst:dv:convertibles-and-the-credit-equity-link:constraints}
\item \textbf{Run} \texttt{dv\_convertible.table()}, \texttt{credit\_sensitivity()}, \texttt{stress(-0.20, 300)}, \texttt{call\_effect()} and
\texttt{fig\_convertible.py}.
\end{tutsteps}
\textbf{What to change next.} Set $p=2$ and recompute the delta at 20; add an investor put at 100 in year three and see the \omterm{def:dv:convertibles-and-the-credit-equity-link:floor}{bond floor} move; price the
call with a hard (unconditional) call from year two.
\end{tutorial}

\section{Build: the convertible pricer}\label{bld:dv:convertibles-and-the-credit-equity-link:convertible}

\begin{build}
\bfield{Purpose} The miniature firm's convertible-bond pricer and risk engine: value, delta, gamma and credit sensitivity of each bond in its
convertible book, under a hazard that may depend on the share.
\bfield{Interface} \texttt{Convertible(face, maturity, coupon, freq, ratio, recovery, call\_start, call\_price, call\_trigger, puts)};
\texttt{power\_hazard(lam0, s0, p, cap)}; \texttt{price\_grid(cb, r, q, vol, hazard, s\_min, s\_max, nx, steps\_per\_year) -> (spots, values)};
\texttt{value\_at}; \texttt{greeks(s0, grid) -> value, delta, gamma}; \texttt{bond\_floor(cb, r, spread)}.
\bfield{Rules} The hazard function is calibrated to the issuer's credit curve before any convertible is priced; conversion, call and put are applied
after every time step; credit sensitivity is reported with the equity delta, and the stress of share and spread together is run daily.
\bfield{Acceptance tests} \texttt{code/firm/convertible/tests/}: without conversion and hazard the price is the bond; with a constant hazard, the
closed form with recovery of face; the value exceeds floor and parity and rises with the share; the call lowers and a put raises the value; the
credit link raises the delta at low share prices.
\bfield{Stretch} The Tsiveriotis--Fernandes split; a hazard calibrated to a CDS curve and to the issuer's options; dividends and call notice
periods; the contingent convertibles of Book 2, chapter 26.
\end{build}

\begin{omsources}
\item K.~Tsiveriotis and C.~Fernandes, ``Valuing convertible bonds with credit risk'', \emph{Journal of Fixed Income} 8(2) (1998) 95--102.
\item E.~Ayache, P.~A.~Forsyth and K.~R.~Vetzal, ``Valuation of convertible bonds with credit risk'', \emph{Journal of Derivatives} 11(1) (2003)
9--29.
\item L.~Andersen and D.~Buffum, ``Calibration and implementation of convertible bond models'', \emph{Journal of Computational Finance} 7 (2003)
1--34.
\item M.~Mitchell, L.~H.~Pedersen and T.~Pulvino, ``Slow moving capital'', \emph{American Economic Review} 97(2) (2007) 215--220.
\item M.~Mitchell and T.~Pulvino, ``Arbitrage crashes and the speed of capital'', \emph{Journal of Financial Economics} 104(3) (2012) 469--490.
\end{omsources}

\section{Exercises}

\begin{exercise}[$\star$]\label{exo:dv:convertibles-and-the-credit-equity-link:1}
A convertible with face 100 and \omterm{def:dv:convertibles-and-the-credit-equity-link:ratio}{conversion ratio} 2.5 trades at 110 with the share at 36. Give the \omterm{def:dv:convertibles-and-the-credit-equity-link:price}{conversion price}, \omterm{def:dv:convertibles-and-the-credit-equity-link:value}{conversion value} and \omterm{def:dv:convertibles-and-the-credit-equity-link:premium}{conversion
premium}.
\end{exercise}

\begin{exercise}[$\star$]\label{exo:dv:convertibles-and-the-credit-equity-link:2}
Why is a convertible worth at least its \omterm{def:dv:convertibles-and-the-credit-equity-link:floor}{bond floor} and at least its \omterm{def:dv:convertibles-and-the-credit-equity-link:value}{conversion value}?
\end{exercise}

\begin{exercise}[$\star$]\label{exo:dv:convertibles-and-the-credit-equity-link:3}
Why does the \omterm{def:dv:convertibles-and-the-credit-equity-link:softcall}{soft call} lower the convertible's value, and where is the effect largest?
\end{exercise}

\begin{exercise}[$\star\star$]\label{exo:dv:convertibles-and-the-credit-equity-link:4}
Explain why the share's drift under the pricing measure is $r-q+\lambda(S)$ in an \omterm{def:dv:convertibles-and-the-credit-equity-link:e2c}{equity-to-credit model}.
\end{exercise}

\begin{exercise}[$\star\star$]\label{exo:dv:convertibles-and-the-credit-equity-link:5}
Why is the convertible's credit sensitivity smaller than its \omterm{def:dv:convertibles-and-the-credit-equity-link:floor}{bond floor}'s?
\end{exercise}

\begin{exercise}[$\star\star$]\label{exo:dv:convertibles-and-the-credit-equity-link:6}
Explain why the equity-to-credit delta is larger than the constant-hazard delta at low share prices, and what that means for a hedge.
\end{exercise}

\begin{exercise}[$\star\star\star$]\label{exo:dv:convertibles-and-the-credit-equity-link:7}
\emph{Coding.} Recompute the stress of \cref{ex:dv:convertibles-and-the-credit-equity-link:stress} with the equity-to-credit exponent $p=2$ (same
spread today). Does the delta-hedged loss grow or shrink?
\end{exercise}

\begin{exercise}[$\star\star\star$]\label{exo:dv:convertibles-and-the-credit-equity-link:8}
\emph{Find the flaw.} ``Our convertible book is delta-hedged and CDS-hedged, so its only exposure is to volatility.''
\end{exercise}

\section{Problem: Equity Down, Credit Wider}

\begin{problem}[{Weekend problem --- the scenario a convertible desk fears}]\label{pb:dv:convertibles-and-the-credit-equity-link:1}
A desk holds 100 million of face of the chapter's convertible (share 30, \omterm{def:dv:convertibles-and-the-credit-equity-link:price}{conversion price} 40, spread 300 basis points), short the equity-to-credit
delta in shares.

\medskip\noindent\textbf{Part I --- The bond.}
\begin{enumerate}
\item Give the convertible's value, \omterm{def:dv:convertibles-and-the-credit-equity-link:floor}{bond floor}, \omterm{def:dv:convertibles-and-the-credit-equity-link:value}{conversion value} and \omterm{def:dv:convertibles-and-the-credit-equity-link:premium}{conversion premium} at 30 under the \omterm{def:dv:convertibles-and-the-credit-equity-link:e2c}{equity-to-credit model}.
\item Give the delta and the number of shares the desk is short.
\item Give the value at a share price of 10 and the hazard there.
\item What does the \omterm{def:dv:convertibles-and-the-credit-equity-link:softcall}{soft call} cost the holder at a share price of 52?
\item Why is the \omterm{def:dv:convertibles-and-the-credit-equity-link:floor}{bond floor} not flat in this model?
\end{enumerate}

\medskip\noindent\textbf{Part II --- The hedges.}
\begin{enumerate}[resume]
\item Give the credit sensitivity per basis point, for the position.
\item Size the CDS protection that offsets it.
\item What does the position earn if the share falls 20\% with the spread unchanged?
\item What does it lose if the spread widens 300 basis points with the share unchanged?
\item What would the constant-hazard model's delta be, and how many fewer shares would the desk be short?
\end{enumerate}

\medskip\noindent\textbf{Part III --- The scenario.}
\begin{enumerate}[resume]
\item The share falls 20\% and the spread at the new price ends 300 basis points wider. Give the convertible's new value and the delta-hedged loss.
\item Give the loss with the constant-hazard delta.
\item Give the result with the CDS hedge as well.
\item Which part of the loss is credit and which is the model's delta?
\item What scenario is missing from this analysis?
\end{enumerate}

\medskip\noindent\textbf{Part IV --- Judgement.}
\begin{enumerate}[resume]
\item Which delta would you trade, and why?
\item How much CDS would you hold, given the jump to default?
\item What limits would you set on the book?
\item State the \emph{named result}: the loss of the delta-hedged convertible-arbitrage position when the issuer's credit spread widens by 300 basis
points while its share falls 20\%.
\item In one sentence: what does a convertible arbitrageur own?
\end{enumerate}
\end{problem}

\section{Interview questions}

\begin{interviewq}[$\star$ \iqroles{trader, researcher}]\label{iq:dv:convertibles-and-the-credit-equity-link:1}
Sketch a convertible's value against the share price. Name the regions.
\end{interviewq}

\begin{interviewq}[$\star\star$ \iqroles{researcher}]\label{iq:dv:convertibles-and-the-credit-equity-link:2}
Write the pricing equation for a convertible with default risk. What are the boundary conditions and constraints?
\end{interviewq}

\begin{interviewq}[$\star\star$ \iqroles{trader}]\label{iq:dv:convertibles-and-the-credit-equity-link:3}
What does a convertible arbitrageur earn, and what are the risks?
\end{interviewq}

\begin{interviewq}[$\star\star$ \iqroles{developer}]\label{iq:dv:convertibles-and-the-credit-equity-link:4}
How would you implement the conversion, call and put features in a finite-difference pricer?
\end{interviewq}

\begin{interviewq}[$\star\star$ \iqroles{risk}]\label{iq:dv:convertibles-and-the-credit-equity-link:5}
Design a stress test for a convertible-arbitrage book.
\end{interviewq}

\begin{interviewq}[$\star\star\star$ \iqroles{trader, risk}]\label{iq:dv:convertibles-and-the-credit-equity-link:6}
Prime brokers cut the book's leverage by half overnight. What happens to the positions and to their prices?
\end{interviewq}
